\documentclass[11pt]{article}

\usepackage[margin=1in]{geometry}
\usepackage{amsmath}
\usepackage{amssymb}
\usepackage{amsthm}
\usepackage{hyperref}
\usepackage{enumitem}

\newtheorem{definition}{Definition}

\title{The Coherence Barcode: \\ A Template-Matching Diagnostic for Ambient-Noise Phase-Velocity Estimation}
\author{FastMSPEC project design document}
\date{2026-08-23}

\begin{document}
\maketitle

\begin{abstract}
This document formalizes a diagnostic and phase-velocity-picking method developed for the
FastMSPEC project: representing a coherence (noise correlation function, NCF) spectrum's
zero-crossings, local maxima, and local minima as a typed sequence of frequency-tagged events
--- a ``barcode'' --- and judging spectral-estimation techniques (single-taper vs.\ multitaper
methods such as FastMspec) by how well their barcode matches a library of physically-motivated
templates, scanned across a plausible range of Love-wave phase velocities. The method is grounded
in, but distinct from, three lines of prior work: Ekstr\"om et al.'s (2009) zero-crossing
phase-velocity construction, Hawkins and Sambridge's (2019) extension of that construction to
peaks and troughs via an analytical Bessel-function result, and Xue and Olugboji's (2025) AkiNet
architecture, whose bounded-corridor reference-curve design independently converges on the same
underlying principle used here. This document is intended as a standalone reference for the
method's design, independent of the conversation history that produced it, so that the
associated notebook implementation can be revisited, audited, or extended without re-deriving
the reasoning from scratch.
\end{abstract}

\tableofcontents

\section{Motivation}

Ambient-noise interferometry estimates the coherence, or noise correlation function (NCF),
between two seismic stations, whose real part is predicted by Aki's (1957) spectral formulation
to follow a Bessel-function form. For Rayleigh waves (and, as a common simplifying
approximation also used throughout this project, for Love waves too --- see
Section~\ref{sec:love-caveat}),
\begin{equation}
  \rho(f, r) \approx J_0\!\left(\frac{2\pi f r}{c(f)}\right),
  \label{eq:aki}
\end{equation}
where $\rho(f,r)$ is the real part of the coherence spectrum, $r$ is the interstation distance,
$f$ is frequency, and $c(f)$ is the (generally frequency-dependent) phase velocity.

Extracting $c(f)$ from an observed, noisy $\rho(f,r)$ is a nonlinear inverse problem. A classical
approach (Ekstr\"om et al., 2009) identifies zero-crossings of the observed spectrum and matches
them against the known zero-crossing structure of Equation~\ref{eq:aki} to construct trial
dispersion curves. The FastMSPEC project (this repository) has, over the course of its
development, repeatedly found that \emph{how well an estimator's coherence spectrum supports
this kind of feature extraction} is itself informative about the estimator's quality --- not just
its variance, but whether its zero-crossings (and, as developed here, its turning points) sit
where a physically plausible earth would put them.

Prior diagnostics built in this project (documented in \texttt{notebooks/03\_fastmspec\_
application.ipynb}, Sections 1c--1f) established two separate findings: (1) single-taper
spectral estimates have much higher variance than multitaper (FastMspec) estimates, and (2)
smoothing single-taper's estimate to match FastMspec's variance does not recover FastMspec's
resolution --- smoothed single-taper's zero-crossings still drift substantially, and the amount
of drift can be predicted from a resolution-bandwidth criterion, $2W \lesssim \Delta f_{\text{zero}}$,
where $\Delta f_{\text{zero}} = c/(2r)$ is the natural zero-crossing spacing implied by the
Aki--Bessel relationship. Both findings, however, only checked \emph{internal self-consistency}
of a candidate spectrum against itself (e.g., how much do a technique's own crossings drift under
smoothing, or how tightly clustered is its own implied zero-crossing spacing). Neither compared a
candidate directly against an actual physical model of what a real earth's coherence spectrum
should look like.

This document formalizes that missing comparison: the \textbf{coherence barcode}.

\section{The Barcode}

\begin{definition}[Barcode]
For any coherence curve $\gamma(f)$ --- real (estimated from data) or theoretical (evaluated from
a model) --- its \emph{barcode} is the ordered sequence of typed, frequency-tagged events:
\begin{itemize}[nosep]
  \item \textbf{Z} (zero-crossing): $f$ such that $\gamma(f) = 0$ with a sign change.
  \item \textbf{M} (maximum): $f$ such that $\gamma$ has a local maximum.
  \item \textbf{N} (minimum): $f$ such that $\gamma$ has a local minimum.
\end{itemize}
\end{definition}

The name and visual framing come directly from how this looks when plotted: each event type
rendered as a distinctly colored vertical mark along the frequency axis produces a pattern
resembling a barcode. This is not merely a visualization convenience --- it is the organizing
metaphor for the whole method: a \emph{good} spectral-estimation technique should produce a
barcode that can be matched, stripe for stripe, against a barcode generated from physically
plausible earth structure.

\subsection{Template barcodes: analytically exact event locations}
\label{sec:template-barcodes}

A \emph{template} is the barcode of Equation~\ref{eq:aki} evaluated at one specific candidate
phase-velocity curve $c(f)$. Because the theoretical curve is smooth and known in closed form
(given $c(f)$), its barcode can be computed exactly, without any of the numerical
sign-change-detection or derivative-estimation machinery needed for noisy real data.

This exactness follows from a classical Bessel-function identity used by Hawkins and Sambridge
(2019, their Equations 1--3), extending the zero-crossing construction of Ekstr\"om et al.\ (2009,
their Equation 2) to peaks and troughs as well:
\begin{itemize}[nosep]
  \item \textbf{Zero-crossings} of $J_0(x)$ occur at $x = j_{0,k}$, the $k$-th zero of $J_0$
    (tabulated; in \texttt{scipy}, \texttt{scipy.special.jn\_zeros(0, N)}).
  \item \textbf{Maxima and minima} of $J_0(x)$ occur where its derivative vanishes. Since
    $\frac{d}{dx} J_0(x) = -J_1(x)$, these occur at $x = j_{1,k}$, the $k$-th zero of $J_1$
    (tabulated via \texttt{scipy.special.jn\_zeros(1, N)}).
\end{itemize}

Given a candidate phase-velocity curve $c(f)$, the Bessel argument is
$x(f) = \frac{2\pi f r}{c(f)}$. Template event frequencies are then obtained by solving
\begin{equation}
  x(f_k) = j_{0,k} \quad \text{(Z events)}, \qquad x(f_k) = j_{1,k} \quad \text{(M/N events)}
\end{equation}
for $f_k$. When $c(f)$ has no closed form (as is the case here, where $c(f)$ is built from a
discretely sampled real dispersion curve --- Section~\ref{sec:refcurve}), this requires a
one-dimensional numerical root-find (e.g., \texttt{scipy.optimize.brentq} against an interpolated
$c(f)$) rather than a closed-form inversion. The key advantage over naive dense-sampling is that
the \emph{target} Bessel-argument value being solved for is exact and tabulated, not approximated
by finite-difference sign-change or derivative detection on a sampled grid --- there is no
discretization error in identifying \emph{which} $x$-values correspond to events, only ordinary
root-finding tolerance in solving for the corresponding $f$.

Maxima and minima alternate strictly with each other and with zero-crossings in a fixed pattern
(since $J_0$ and $J_1$'s zeros interlace); this provides a useful internal consistency check when
implementing the solver (Section~\ref{sec:verification}).

\subsection{Candidate barcodes: noise-contaminated, need reliability filtering}

A \emph{candidate} barcode is extracted from a real spectral estimate --- FastMspec, single-taper
(raw or smoothed), Mspec, or MspecBestK. Unlike templates, candidates have no closed form; their
barcodes must be extracted numerically (sign-change detection for Z, local extrema detection for
M/N) on the sampled, noisy estimate itself.

Real spectra contain spurious noise-driven events alongside genuine ones. Prior work in this
project (Notebook 3, Section 1f) established that these cannot be separated by a single absolute
threshold shared across techniques: a low-variance estimate (FastMspec) and a high-variance one
(raw single-taper) have incomparable absolute amplitude-swing scales --- a threshold calibrated
from FastMspec's swing distribution would trivially pass 100\% of single-taper's noise-driven
events, since single-taper's noise alone produces larger local swings than FastMspec's entire
signal.

The resolution, carried forward here: each method's events are filtered to ``reliable'' using
that method's \emph{own} relative threshold --- e.g., the top 50\% by local amplitude swing
(maximum minus minimum in a small window straddling the event) within that method's own
distribution, not an absolute value shared across techniques. This is applied uniformly to Z, M,
and N events (turning-point prominence generalizes the same construction used for zero-crossing
swing).

\section{Matching and Scoring}
\label{sec:scoring}

Given a candidate barcode and one template barcode, a \emph{match score} quantifies how well they
agree. Four design elements, discussed in turn below: (1) a same-type, locally-scaled matching
tolerance; (2) combined precision and recall, not recall alone; (3) hard rejection of
physically-implausible matches, mirroring AkiNet's architectural (not merely penalized) bounding
corridor; and (4) frequency-dependent trust weighting, following an explicit argument in Hawkins
and Sambridge (2019) about where a reference earth model is most and least reliable.

\subsection{Matching tolerance}

A candidate event matches a template event of the same type if the two frequencies are within a
tolerance. This tolerance is scaled to the \emph{template's own local event spacing} at that
point, not a fixed frequency value --- consistent with the drift-as-a-percentage-of-local-spacing
methodology already validated in Notebook 3 Section 1f. A fixed absolute tolerance would be too
loose where template events are closely spaced (high frequency) and too tight where they are
widely spaced (low frequency).

\subsection{Precision and recall, not recall alone}

A naive score --- the fraction of template events that have a matching candidate event (recall)
--- is vulnerable to a specific failure mode: a very dense, noisy candidate can trivially
``match'' almost any template purely by having many more events than the template does. Concretely,
in this project's own data, raw single-taper produces roughly 1{,}135 zero-crossings over a band
where FastMspec produces 126; recall-only scoring would reward exactly the noise-dominated
behavior the diagnostic is meant to penalize. The score therefore combines precision (fraction of
candidate events that correspond to a real template event) and recall (fraction of template
events that are matched) into a combined measure (e.g., their harmonic mean, analogous to an
F1 score).

\subsection{Hard corridor rejection}
\label{sec:corridor-rejection}

AkiNet (Xue and Olugboji, 2025, Section 3.3.1) bounds its inverted phase velocity to
$c^{\text{ivt}}(f) = W\hat{z}(f) + b(f)$, where $b(f)$ is a reference curve and $W$ is a diagonal
scaling matrix defining a fixed perturbation corridor (they use $\pm 0.8$~km/s) around it. This is
an \emph{architectural} constraint: the network's output cannot leave the corridor regardless of
what the data appear to want, not a soft penalty term traded off against data fit in a loss
function.

This method mirrors that design choice directly, rather than relying only on soft
precision/recall scoring: a candidate event whose \emph{implied} local phase velocity (via the
Ekstr\"om et al.\ 2009-style relation $c = \omega r / z$, evaluated at that event's own frequency
and nearest Bessel-zero index) falls outside the \emph{entire} template corridor --- not merely
``does not match the single best-fitting template,'' but ``is not explainable by any corridor
member at all'' --- is rejected as a candidate match, rather than merely counted as one more
unmatched (but otherwise unremarkable) event. This distinguishes ``noisy, but plausible'' events
from ``implying an unphysical velocity outright,'' which a flat precision/recall score alone
cannot.

\subsection{Frequency-dependent trust}
\label{sec:freq-weighting}

Hawkins and Sambridge (2019) argue directly for why a reference earth model should be trusted
more at some frequencies than others: ``most of the deviation between a prior Earth model and a
new regional study will be confined to the near surface or higher frequency'' (Discussion
section). Long-period (low-frequency) surface waves sample deeper, larger-scale, more spatially
averaged structure, which tends to vary less between a generic reference model and any specific
region; short-period (high-frequency) waves sample shallow, heterogeneous structure where
regional deviation from a reference model is expected and legitimate.

This method therefore weights mismatches (missed template events, spurious candidate events, or
corridor-rejected events, Section~\ref{sec:corridor-rejection}) \emph{more heavily} at the
low-frequency end of the analysis band than at the high-frequency end. A technique that fails to
reproduce the barcode structure where the reference model is most trustworthy is more damning
than one that struggles only where legitimate regional deviation is expected. This replaces a
naive frequency-uniform scoring scheme.

\subsection{Diagnostic value of a null result}

AkiNet's own discussion frames a specific failure mode --- the network's output defaulting to the
bare reference curve $b(f)$ (i.e., zero perturbation) --- as a \emph{useful, diagnostic} failure
signature, not merely a bad outcome: it occurs when the input data are too poor to inform the fit
at all, and its predictability makes such cases easy to flag for quality control (Xue and
Olugboji, 2025, Section 4.4 and Discussion).

This method surfaces the analogous case explicitly: if a candidate's best-matching template turns
out to be (or lies very close to) $\delta = 0$ (the bare reference curve, no perturbation), this
is reported alongside the match score rather than silently absorbed into a single scalar. Combined
with a poor absolute match score, it indicates the candidate carries too little real signal to
usefully constrain anything, rather than being read as a confident ``matches the reference well''
result.

\section{The Template Family}
\label{sec:refcurve}

\subsection{Reference curve: \texttt{SDISPL.ASC}}

Templates require a base reference curve $c_{\text{ref}}(f)$ to build the perturbed family around
(Section~\ref{sec:template-family}). A rough, hand-placed reference curve already existed
elsewhere in this project (Notebook 3, Section 4, used for the \texttt{seislib} comparison), but
its provenance was undocumented and it was explicitly flagged in this project's own notes as ``NOT
Sayan's exact curve.''

A real, physically computed reference curve was located during this design process, on the lab
network-attached storage (\texttt{repovibranium}), at
\texttt{madagascar\_data/pre\_computed\_files/SDISPL.ASC}. This is a standard Herrmann
Computer-Programs-in-Seismology-format ASCII dispersion-curve file (columns: \texttt{LMODE},
\texttt{NFREQ}, \texttt{PERIOD(S)}, \texttt{FREQUENCY(Hz)}, \texttt{C(KM/S)}), containing a
multi-mode Love-wave dispersion computation: the fundamental mode (\texttt{LMODE=0}) spans periods
2--2048~s (velocities 1.22--4.83~km/s), with higher modes (\texttt{LMODE=1--8}) each contributing
starting at progressively shorter periods. This is the same file used by Sayan Swar's own
\texttt{phasevel\_compute\_slide13and14.ipynb} (also on \texttt{repovibranium}, in
\texttt{madagascar\_codes/python/}) to build the reference curve shown in his project report's
Figure~6, via exactly the same \texttt{LMODE==0} filter used here.

Restricted to this project's actual analysis band (0.02--0.3~Hz), the fundamental mode shows
normal dispersion (velocity decreasing with increasing frequency): 4.49~km/s at 0.02~Hz down to
1.53~km/s at 0.3~Hz. This range brackets a much cruder empirical estimate derived independently
elsewhere in this project (Notebook 3, Section 1f: $\sim$2.11~km/s from raw zero-crossing
spacing) sensibly, which is itself a useful cross-check on both.

\subsection{Building the template family}
\label{sec:template-family}

A single reference curve is not, by itself, ``a range of plausible phase velocities.'' Following
AkiNet's own design (Section~\ref{sec:corridor-rejection}), the template family is built as an
\emph{additive} perturbation corridor around the reference curve:
\begin{equation}
  c_{\text{template}}(f) = c_{\text{ref}}(f) + \delta, \qquad \delta \in [-0.8, 0.8] \text{ km/s},
\end{equation}
stepped finely (e.g.\ 0.05~km/s steps, giving roughly 33 templates across the corridor). The
corridor width, $\pm 0.8$~km/s, is adopted directly from AkiNet's own tuned value --- chosen by
its authors, by experimentation, as ``generous enough to span all plausible velocities required by
the data'' on the same problem domain (real Love-wave ambient noise) --- rather than re-derived
independently.

An earlier version of this design considered \emph{multiplicative} scaling of the reference curve
instead of additive perturbation; this was revised after reviewing AkiNet's own architecture,
which is explicitly additive, not multiplicative.

\subsection{The Love-wave Bessel-model simplification}
\label{sec:love-caveat}

Equation~\ref{eq:aki} (pure $J_0$) is, strictly, the Rayleigh-wave form. AkiNet (Xue and Olugboji,
2025, Equation 3, citing Nishida et al., 2024) gives the more physically correct Love-wave form as
a combination of the zeroth- and second-order Bessel functions,
\begin{equation}
  \rho(f, r) \approx \frac{1}{2}\left[J_0\!\left(\frac{2\pi f r}{c(f)}\right) - J_2\!\left(\frac{2\pi f r}{c(f)}\right)\right].
\end{equation}
This method uses pure $J_0$ for its first version, matching every prior use of the Bessel model
elsewhere in this project (Notebook 2's illustrative example, \texttt{nb3\_helpers.bessel\_fit\_
quality}, and Notebook 3 Section 1d's synthetic leakage test), and matching Hawkins and
Sambridge's (2019) own practical simplification --- their paper uses pure $J_0$ for both Love and
Rayleigh data throughout, despite studying both wave types.

The trade-off is explicit: pure $J_0$ keeps the exact, tabulated-zero property that makes the
analytical approach in Section~\ref{sec:template-barcodes} valuable in the first place. The
combined $J_0 - J_2$ form does not have simply tabulated zeros (it is a difference of two
different Bessel functions, not one), and would require numerical root-finding for template
events too, losing most of the practical advantage of the analytical approach for this first
version. This is recorded here as a known, deliberate simplification, not an oversight --- a
natural candidate for a future revision of this method.

\section{Relationship to Existing Tools}

This method is a genuine synthesis, not a restatement of any single existing tool in this project
or in the literature reviewed while designing it:

\begin{itemize}
  \item \textbf{\texttt{nb3\_helpers.bessel\_fit\_quality}} (this project) grid-searches a
    constant-velocity trial $c$ only (never a frequency-dependent $c(f)$), and fits the
    \emph{full complex amplitude} of the observed coherence against the Bessel prediction via RMS
    residual. This project's own documentation already flags the resulting metric as unreliable,
    since it is ``dominated by overall coherence strength, not Bessel-shape quality.'' A
    crossing/turning-point barcode is structurally immune to this specific failure mode: event
    \emph{locations} do not depend on the overall amplitude scale the way a full RMS-amplitude
    fit does.
  \item \textbf{Notebook 3, Section 1f} (this project) checks only internal self-consistency of a
    candidate's own crossing spacing against itself; it never compares a candidate against an
    actual external physical template. This method's template library supplies exactly that
    missing external reference.
  \item \textbf{\texttt{seislib.extract\_dispcurve}} (Magrini et al., 2022) uses zero-crossings
    only (no turning points) and a single fixed reference curve, used as a mode-tracking
    guide/prior with a tolerance band around the reference velocity at each step --- confirmed by
    reading its source (\texttt{\_an\_processing\_seislib\_functions.py}, obtained from the same
    lab NAS location as the reference curve used here). This is a sequential-tracking mechanism,
    not an explicit multi-hypothesis scan-and-score across a template library. AkiNet's own paper
    (Xue and Olugboji, 2025) identifies ``cycle skip'' --- committing early to the wrong phase
    oscillation --- as exactly this class of method's characteristic failure mode; a
    scan-everything, score-everything approach is structurally more resistant to it, since it
    never commits to a single crossing sequence before evaluating alternatives.
  \item \textbf{AkiNet} (Xue and Olugboji, 2025) solves the same underlying inverse problem via
    continuous physics-informed neural-network optimization rather than a discrete scan over
    templates. Its bounded-corridor architecture is philosophically identical to the template
    family used here (Section~\ref{sec:template-family}) --- both encode ``a plausible
    neighborhood around a physically motivated reference,'' arrived at independently via two very
    different mechanisms.
  \item \textbf{Ekstr\"om et al.\ (2009)} already constructs multiple trial phase-velocity curves
    from observed zero-crossings (indexed by an integer ambiguity $m$ in their Equation~2) and
    selects whichever is closest to a reference curve --- the same core ``generate candidates,
    compare against a reference, pick the best'' idea this method formalizes and extends, using
    an explicit multi-template precision/recall score (Section~\ref{sec:scoring}) rather than
    single-curve closest-match selection, and three event types instead of one.
  \item \textbf{Hawkins and Sambridge (2019)} extend that same zero-crossing construction to
    peaks and troughs via the $J_1$-zero identity (Section~\ref{sec:template-barcodes}) --- the
    direct analytical basis this method's template-event generation relies on --- within a
    fundamentally different overall method (an adjoint-based nonlinear waveform-fitting
    inversion of a full path-average earth model, not a discrete barcode-matching scan).
\end{itemize}

The barcode/template-matching framing itself --- three typed event streams, matched via combined
precision, recall, corridor rejection, and frequency-weighted trust against a scanned template
library --- is original to this project, not a restatement of any of the above.

\section{Scope and Open Items}

\begin{itemize}[nosep]
  \item \textbf{Deliverable.} Both a reliability metric (for judging FastMspec against
    single-taper and other techniques) and a phase-velocity pick (the best-matching template's
    $c(f)$) --- the latter a from-scratch alternative to \texttt{seislib.extract\_dispcurve}, not
    a replacement for Notebook 3's existing Section~4, which remains as independent, unmodified
    work.
  \item \textbf{Data scope.} SA53/SA58 (Love-wave, transverse component) only for the first
    version. MTAN/RUNG --- the pair where \texttt{seislib} previously failed to converge in this
    project, and the same pair AkiNet's own paper uses as a stress-test example --- is an obvious
    and motivated next test, deferred to a later round once this method is validated on the
    cleaner pair.
  \item \textbf{Event weighting.} Z, M, and N event types are weighted equally in the match score
    for this first version; revisiting this (e.g. weighting crossings more heavily, since they
    tend to be more sharply localized than turning points) is a plausible future refinement.
  \item \textbf{Love-wave Bessel model.} Pure $J_0$ for this first version
    (Section~\ref{sec:love-caveat}); the more physically correct $J_0 - J_2$ combined form is a
    documented, deliberate simplification, not an oversight.
\end{itemize}

\section{Implementation Notes}
\label{sec:verification}

The companion notebook implementation should verify, before any results are trusted or reported:
\begin{enumerate}
  \item \texttt{SDISPL.ASC} parses identically to Sayan Swar's own approach (\texttt{LMODE==0}
    filter, \texttt{[FREQUENCY\_Hz, C\_KM\_S]} columns).
  \item The template library is a sensible, non-degenerate fan of curves when plotted
    ($c_{\text{ref}}(f)$ and several $c_{\text{ref}}(f) + \delta$ members over 0.02--0.3~Hz).
  \item The analytical event solver (Section~\ref{sec:template-barcodes}) is correct: root-found
    frequencies should zero out $J_0(x(f))$ (or $J_1(x(f))$ for extrema) to near machine
    precision, and should agree with a brute-force dense-sampling/sign-change check on the same
    template curve; Z/M/N events should alternate in the expected order.
  \item FastMspec's best-match score should exceed single-taper's (raw and roughness-matched
    smoothed variants) on real data --- the expected result given everything established in
    Notebook 3 Sections 1c--1f. A result to the contrary should be treated as a bug to find, not a
    finding to report.
  \item The picked best-fit velocity should be checked against both the raw reference curve's own
    value in the same band and the cruder empirical estimate from Section 1f, as a sanity range
    check.
\end{enumerate}

\begin{thebibliography}{9}

\bibitem{aki1957}
Aki, K. (1957). Space and time spectra of stationary stochastic waves, with special reference to
microtremors. \textit{Bulletin of the Earthquake Research Institute}, 35, 415--456.

\bibitem{ekstrom2009}
Ekstr\"om, G., Abers, G.\ A., \& Webb, S.\ C.\ (2009). Determination of surface-wave phase
velocities across USArray from noise and Aki's spectral formulation. \textit{Geophysical Research
Letters}, 36(18), L18301. \url{https://doi.org/10.1029/2009GL039131}

\bibitem{ekstrom2014}
Ekstr\"om, G.\ (2014). Love and Rayleigh phase-velocity maps, 5--40~s, of the western and central
USA from USArray data. \textit{Earth and Planetary Science Letters}, 402, 42--49.
\url{https://doi.org/10.1016/j.epsl.2013.11.022}

\bibitem{ekstrom2017}
Ekstr\"om, G.\ (2017). Short-period surface-wave phase velocities across the conterminous United
States. \textit{Physics of the Earth and Planetary Interiors}, 270, 168--175.
\url{https://doi.org/10.1016/j.pepi.2017.07.010}

\bibitem{hawkins2019}
Hawkins, R., \& Sambridge, M.\ (2019). An adjoint technique for estimation of interstation phase
and group dispersion from ambient noise cross correlations. \textit{Bulletin of the Seismological
Society of America}, 109(5), 1716--1728. \url{https://doi.org/10.1785/0120190060}

\bibitem{magrini2022}
Magrini, F., Lauro, S., K\"astle, E., \& Boschi, L.\ (2022). Surface-wave tomography using
SeisLib: a Python package for multiscale seismic imaging. \textit{Geophysical Journal
International}, 231(2), 1011--1030. \url{https://doi.org/10.1093/gji/ggac236}

\bibitem{xue2025}
Xue, S., \& Olugboji, T.\ (2025). AkiNet: A physics-informed AI for wave extraction from noise.
\textit{Journal of Geophysical Research: Machine Learning and Computation}, 2(4), e2025JH000932.
\url{https://doi.org/10.1029/2025JH000932}

\bibitem{olugboji2022}
Olugboji, T., \& Xue, S.\ (2022). A short-period surface-wave dispersion dataset for model
assessment of Africa's crust: ADAMA. \textit{Seismological Research Letters}, 93(3), 1943--1959.
\url{https://doi.org/10.1785/0220210355}

\end{thebibliography}

\end{document}
